% !TEX root = ../main.tex
\clearpage

We now turn to Zeeprio and begin with coercion resistance. The key observation we call non-expansivity: Zeeprio does not reveal any information beyond what is already revealed by Eperio. More precisely, any information available to a coercer from a Zeeprio transcript can be derived from an Eperio transcript, possibly with weaker assurance. Consequently, Zeeprio cannot be less coercion resistant than Eperio: any privacy property that holds for Eperio continues to hold when augmented with Zeeprio’s additional verification artifacts. 

In particular, if non-expansivity holds, Zeeprio inherits at least the level of privacy established for Eperio by Basin \etal~\cite{BDGR21}, namely receipt-freeness under their table-based definition. Moreover, Zeeprio may also inherit coercion resistance in the stronger sense of Küsters \etal~\cite{KTV10c}, provided one accepts the conjecture that replacing Scantegrity II’s switchboard with Eperio does not reduce privacy in any meaningful way.

We formalize this in the following lemma:

\begin{lemma}[Non-Expansive Privacy under Transcript Augmentation]
Let $\Pi$ be a protocol and let $\Pi^{+}$ be an augmentation that, on the same underlying execution, produces a coercer view
$T^{+} := (T,Z)$, where $T$ is the coercer view in $\Pi$ and $Z$ is an additional transcript component.
Assume that for every probabilistic polynomial-time distinguisher $\mathcal{D}$ and every security parameter $\lambda$,
there exists a probabilistic polynomial-time algorithm $\mathsf{Sim}$ such that
\[
\Bigl|\Pr[\mathcal{D}(T,Z)=1] - \Pr[\mathcal{D}(T,\mathsf{Sim}(T))=1]\Bigr|
\]
is negligible as a function of $\lambda$.
\end{lemma}

Under this lemma, $\Pi^{+}$ is not less private than $\Pi$ for any privacy notion whose experiment depends on the protocol only through
the coercer view. In particular, for any such privacy experiment and any probabilistic polynomial-time adversary
$\mathcal{A}$ interacting with $\Pi^{+}$, there exists a probabilistic polynomial-time adversary
$\mathcal{B}$ interacting with $\Pi$ such that the difference between the success probabilities of
$\mathcal{A}$ and $\mathcal{B}$ in their respective experiments is negligible in $\lambda$.

 In the special case where $Z$ is efficiently computable from $T$ (and public parameters), the simulator
$\mathsf{Sim}$ may be chosen to simply compute $Z$ from $T$. In other words, $\mathsf{Sim}$ is pure post-processing. This will be the dominant case for Zeeperio.

To reason about the non-expansive privacy of Zeeprio, we separate two concerns: (i) the \emph{cryptographic} claim that the proof a polynomial vanishes is zero-knowledge and proves nothing beyond the statement that the polynomial vanishes; and (i) the \emph{semantic} claim that the constraints implied by Zeeperio are non-expansive over an Eperio transcript.

Starting with the \emph{cryptographic} claim, we directly assume the PIOP used by Zeeprio satisfies the standard zero-knowledge property: its proof transcript reveals no information beyond the fact that the asserted algebraic relations are true. Concretely, as formalized for PLONK-style polynomial IOPs, there exists a simulator that, given only the statement that the constraints are satisfied, produces a transcript that is indistinguishable from an honestly generated one. This simulator construction is explicit in the PIOP literature~\cite{Sef24,Lip25} and does not depend on the particular semantics of the underlying constraints, so we do not reprove anything regarding this.

We do however add two important implementation details that allow the simulation to go through. The first is that the polynomial commitment scheme needs to be hiding. When using KZG commitments~\cite{KZG10}, it is sufficient to use their second construction informally referred to as the Pedersen-variant of KZG. The second is that the PIOP proof will open committed polynomials at random challenge points. If the adversary has a guess as to the specific polynomial under commitment, it can use this random challenge to confirm its guess. To prevent this, polynomials should incorporate extra random points (off-domain) at interpolation time, where the number of points exceed the number of points that will be eventually opened on the given polynomial. This lacks a standard name in the literature and we propose calling it the MBKM heuristic~\cite{MBKM19} after it was popularized in the zk-SNARK system Sonic.

We spend the rest of the proof dealing with the \emph{semantic} claim: the algebraic relations asserted by Zeeprio’s constraint system encode no more information than what is already revealed by Eperio’s verification logic.


\clearpage
